% GENERATED FILE. Edit canonical lessons, then run npm run generate:books.
% canonical-source-hash: fnv1a64:66442d174f35a1c2
% canonical-lessons: LA-C188-architectus, LA-C188-textor, LA-C188-figulus, LA-C188-vilicus, LA-C188-colonus

\chapter{Architect, Weaver, Potter, Farm manager, Tenant farmer}
\label{ch:la-architect-weaver-potter-farm-manager-tenant-farmer}
\hlchaptermodality{188}

\begin{chapteropening}
\textbf{By the end of this chapter:} \emph{I can say an architect, a weaver, a potter, a farm manager and a tenant farmer.}

Complete the last lesson of chapter 188: I can say an architect, a weaver, a potter, a farm manager and a tenant farmer.
\end{chapteropening}

\section[architectus, architectī]{architectus, architectī — an architect}

\label{lesson:LA-C188-architectus}



\begin{quote}
Before the new one: say the Latin for freedom, then the Latin for glory.
\end{quote}

\subsection*{You'll want to know: architectus, architectī}
\textbf{architectus, architectī} — \textquotedblleft{}an architect\textquotedblright{}.

\begin{cousinweb}
\textbf{architectus} was borrowed from Greek \emph{arkhitektōn}, \textquotedblleft{}a master builder\textquotedblright{} (Lewis \& Short), from \emph{arkhi-}, \textquotedblleft{}chief\textquotedblright{}, and \emph{tektōn}, \textquotedblleft{}builder\textquotedblright{} (EtymDB 2.1).
\end{cousinweb}

\begin{grammarlens}[title={What you've built}]
One more word for this chapter.
\end{grammarlens}

\subsection*{Your turn}
\begin{itemize}
\raggedright
  \item \emph{Say it:} \emph{architectus, architectī}
  \item \emph{Say it:} \emph{architectus, architectī}, once more
  \item \emph{Say it:} say \emph{architectus, architectī}
  \item \emph{Recall:} say the Latin for freedom, then the Latin for glory, then say \emph{architectus, architectī} again
\end{itemize}

\begin{tcolorbox}[breakable,colback=teal!4,colframe=teal!35!black,title={Before you move on}]
What does architectus, architectī mean? (\textquotedblleft{}An architect\textquotedblright{}.) Say it once more.
\end{tcolorbox}
\section[textor, textōris]{textor, textōris — a weaver}

\label{lesson:LA-C188-textor}



\begin{quote}
Before the new one: say the Latin for glory, then the Latin for an architect.
\end{quote}

\subsection*{You'll want to know: textor, textōris}
\textbf{textor, textōris} — \textquotedblleft{}a weaver\textquotedblright{}.

\begin{cousinweb}
\textbf{textor} is built on \emph{texō}, \textquotedblleft{}to weave\textquotedblright{} (Lewis \& Short), with \emph{-tor}, \textquotedblleft{}one who does it\textquotedblright{}. \emph{texō} is usually traced to Proto-Indo-European \emph{*teks-}, \textquotedblleft{}to weave, to fashion\textquotedblright{}; English \emph{text} and \emph{textile} come from it.
\end{cousinweb}

\begin{grammarlens}[title={What you've built}]
One more word for this chapter.
\end{grammarlens}

\subsection*{Your turn}
\begin{itemize}
\raggedright
  \item \emph{Say it:} \emph{textor, textōris}
  \item \emph{Say it:} \emph{textor, textōris}, once more
  \item \emph{Say it:} say \emph{textor, textōris}
  \item \emph{Recall:} say the Latin for glory, then the Latin for an architect, then say \emph{textor, textōris} again
\end{itemize}

\begin{tcolorbox}[breakable,colback=teal!4,colframe=teal!35!black,title={Before you move on}]
What does textor, textōris mean? (\textquotedblleft{}A weaver\textquotedblright{}.) Say it once more.
\end{tcolorbox}
\section[figulus, figulī]{figulus, figulī — a potter}

\label{lesson:LA-C188-figulus}



\begin{quote}
Before the new one: say the Latin for an architect, then the Latin for a weaver.
\end{quote}

\subsection*{You'll want to know: figulus, figulī}
\textbf{figulus, figulī} — \textquotedblleft{}a potter\textquotedblright{}.

\begin{cousinweb}
\textbf{figulus} is built on \emph{fingō}, \textquotedblleft{}to shape, to mould\textquotedblright{} (Lewis \& Short), on its root without the n: the moulder of clay.
\end{cousinweb}

\begin{grammarlens}[title={What you've built}]
One more word for this chapter.
\end{grammarlens}

\subsection*{Your turn}
\begin{itemize}
\raggedright
  \item \emph{Say it:} \emph{figulus, figulī}
  \item \emph{Say it:} \emph{figulus, figulī}, once more
  \item \emph{Say it:} say \emph{figulus, figulī}
  \item \emph{Recall:} say the Latin for an architect, then the Latin for a weaver, then say \emph{figulus, figulī} again
\end{itemize}

\begin{tcolorbox}[breakable,colback=teal!4,colframe=teal!35!black,title={Before you move on}]
What does figulus, figulī mean? (\textquotedblleft{}A potter\textquotedblright{}.) Say it once more.
\end{tcolorbox}
\section[vīlicus, vīlicī]{vīlicus, vīlicī — a farm manager}

\label{lesson:LA-C188-vilicus}



\begin{quote}
Before the new one: say the Latin for a weaver, then the Latin for a potter.
\end{quote}

\subsection*{You'll want to know: vīlicus, vīlicī}
\textbf{vīlicus, vīlicī} — \textquotedblleft{}a farm manager\textquotedblright{}.

\begin{cousinweb}
\textbf{vīlicus} is built on \emph{vīlla}, \textquotedblleft{}a farmhouse\textquotedblright{} (Lewis \& Short): the man in charge of the farm.
\end{cousinweb}

\begin{grammarlens}[title={What you've built}]
One more word for this chapter.
\end{grammarlens}

\subsection*{Your turn}
\begin{itemize}
\raggedright
  \item \emph{Say it:} \emph{vīlicus, vīlicī}
  \item \emph{Say it:} \emph{vīlicus, vīlicī}, once more
  \item \emph{Say it:} say \emph{vīlicus, vīlicī}
  \item \emph{Recall:} say the Latin for a weaver, then the Latin for a potter, then say \emph{vīlicus, vīlicī} again
\end{itemize}

\begin{tcolorbox}[breakable,colback=teal!4,colframe=teal!35!black,title={Before you move on}]
What does vīlicus, vīlicī mean? (\textquotedblleft{}A farm manager\textquotedblright{}.) Say it once more.
\end{tcolorbox}
\section[colōnus, colōnī]{colōnus, colōnī — a tenant farmer}

\label{lesson:LA-C188-colonus}



\begin{quote}
Before the new one: say the Latin for a potter, then the Latin for a farm manager.
\end{quote}

\subsection*{You'll want to know: colōnus, colōnī}
\textbf{colōnus, colōnī} — \textquotedblleft{}a tenant farmer\textquotedblright{}.

\begin{cousinweb}
\textbf{colōnus} is built on \emph{colō}, \textquotedblleft{}to till, to tend\textquotedblright{} (Lewis \& Short): one who tills the land. \emph{colō} is usually traced to Proto-Indo-European \emph{*kwel-}, \textquotedblleft{}to turn, to move about\textquotedblright{} (EtymDB 2.1). English \emph{colony} comes from the same family.
\end{cousinweb}

\begin{grammarlens}[title={What you've built}]
One more word for this chapter.
\end{grammarlens}

\subsection*{Your turn}
\begin{itemize}
\raggedright
  \item \emph{Say it:} \emph{colōnus, colōnī}
  \item \emph{Say it:} \emph{colōnus, colōnī}, once more
  \item \emph{Say it:} say \emph{colōnus, colōnī}
  \item \emph{Recall:} say the Latin for a potter, then the Latin for a farm manager, then say \emph{colōnus, colōnī} again
\end{itemize}

\begin{tcolorbox}[breakable,colback=teal!4,colframe=teal!35!black,title={Before you move on}]
What does colōnus, colōnī mean? (\textquotedblleft{}A tenant farmer\textquotedblright{}.) Say it once more.
\end{tcolorbox}
